% Exercises the capture of cross-references and citations in the *body* of a
% result, not only in its proof and its optional title. The statements option
% governs it and is on by default, so it is deliberately not given here; the
% options fixture covers statements=false. Checks that: a \ref and a \cite in a
% statement become dependencies of that result; an optional title still
% contributes its own edge (and is not counted twice); the capture stops at
% \end, so a reference in ordinary text between results adds no edge; and a
% proof following a result still works.
% citelabel=key keeps the fixture bibliography-style independent, and the
% manual \bibitem keeps it bibtex-free.
\documentclass{article}
\usepackage{amsthm}
\usepackage[cite,citelabel=key]{proofgraph}

\newtheorem{theorem}{Theorem}
\newtheorem{lemma}{Lemma}

\begin{document}

\begin{lemma}\label{l1}One.\end{lemma}

\begin{theorem}\label{t1}
  Under the hypotheses of \ref{l1}, and of \cite[Thm~2]{bk}, something holds.
\end{theorem}

By \ref{l1}, this sentence is ordinary text and adds no edge.

\begin{theorem}[(refines \ref{t1})]\label{t2}
  Something better still, using \cite{bk}.
\end{theorem}
\begin{proof}By \ref{l1}.\end{proof}

\begin{thebibliography}{9}\bibitem{bk}A book.\end{thebibliography}
\end{document}
